\section{A persistent block-structure entropy deficit}\label{sec:blocks}

We now bound the entropy of the particular terminal ensembles retained
in Lemma~\ref{lem:flow}. Fix its cutoff $R$, dimension subsequence, and
uniform microstate ensembles $A^{(d)}$. The parameters $p,h,t$ remain
fixed as in Section~\ref{sec:signal}. For a decomposition of the identity
into $p$ orthogonal projections $P=(P_l)$, write
\[
 \D_P D=\left(\sum_l P_lD_iP_l\right)_i,
 \qquad D_{[p]}=(D_1,\ldots,D_p).
\]
Use the midpoint spectral partition of the last coordinate of $A^{(d)}$
to obtain $P$. For every tuple in its support and all sufficiently large
dimensions,
\begin{equation}\label{eq:matrixblocks}
 \rank(P_l)\leq\frac{2d}{p},\qquad
 \|\D_P A^{(d)}_{[p]}\|_{2,d}^2>p/2.
\end{equation}
To prove this uniformly, first consider any deterministic sequence of
tuples at cutoff $R$ whose joint moments converge to those of $X$.
Write $I_l$ for the midpoint intervals, with any fixed convention at
their endpoints, and let $\nu_d$ be the empirical spectral measure of
$A_n^{(d)}$. These measures converge weakly to the law of $X_n$.
Choose a union $N$ of small closed neighborhoods of the midpoint
thresholds, disjoint from $\sigma(X_n)$. On a common compact interval containing
$[-R,R]$ and $\sigma(X_n)$, choose continuous functions
$0\leq f_l\leq1$ agreeing with $1_{I_l}$ outside $N$. Then
$f_l(X_n)=e_l$, and, with $F_l^{(d)}=f_l(A_n^{(d)})$,
\[
 \|P_l-F_l^{(d)}\|_{2,d}^2\leq\nu_d(N)\longrightarrow0.
\]
Here weak convergence gives $\nu_d(N)\to0$ because $N$ is closed and
has zero limiting mass. In particular, small spectral outliers and
eigenvalues exactly at a threshold are allowed.

The cutoff controls multiplication by every other coordinate:
\begin{equation}\label{eq:block-outlier}
 \|P_lA_i^{(d)}P_l-F_l^{(d)}A_i^{(d)}F_l^{(d)}\|_{2,d}
 \leq2R\|P_l-F_l^{(d)}\|_{2,d}.
\end{equation}
The two compressed norms are at most $R$, so their squared norms also
have difference tending to zero. Uniform polynomial approximation of
$f_l$ and joint moment convergence now imply
\[
 \frac{\rank(P_l)}d\longrightarrow\tau(e_l),\qquad
 \|P_lA_i^{(d)}P_l\|_{2,d}^2
 \longrightarrow\|e_lX_i e_l\|_2^2.
\]
The block summands are orthogonal in Hilbert--Schmidt inner product;
hence their squared norms add to
$\|\D_P A^{(d)}_{[p]}\|_{2,d}^2$. The strict margins in
\eqref{eq:smallblocks} give \eqref{eq:matrixblocks}.
If uniformity over the selected supports failed, a violating tuple in
each of a subsequence of dimensions would itself be a bounded
moment-convergent sequence, contradicting this argument. Closure points
of the microstate sets satisfy the same moment bounds with weak
inequalities, whose errors still tend to zero.

Put $v=1+h+t$, and let $\mathcal E_d$ be the following event for a matrix
tuple $D$:
\begin{equation}\label{eq:event}
 \begin{gathered}
 \|D_{[p]}\|_{2,d}\leq4\sqrt{vp},\qquad
 \text{there exists a decomposition }P\text{ with}\\
 \rank(P_l)\leq2d/p,\quad
 \|\D_P D_{[p]}\|_{2,d}\geq\sqrt p/4.
 \end{gathered}
\end{equation}
This is a deterministic Borel event. For each admissible rank list, the
space of projection decompositions is a compact unitary orbit. The
maximum of $\|\D_P D_{[p]}\|_{2,d}$ over that orbit is a continuous
function of $D$: every compression is contractive, so the maximum is
even Lipschitz in the tuple $2$-norm. The event is therefore a finite
union of closed threshold sets intersected with a closed norm ball.
Its definition does not depend on $A^{(d)}$; the spectral decomposition
of $A_n^{(d)}$ will serve only as a witness for its probability under
the evolved ensemble.

\begin{lemma}\label{lem:eventprob}
Let $\mu_d$ be the law of $A^{(d)}+\sqrt t G^{(d)}$, and let $\gamma_d$
be the law of independent centered GUE coordinates of variances
$v_i(h+t)=h+t+c_i$. For all sufficiently large $d$,
\[
 \mu_d(\mathcal E_d)\geq\tfrac12,
 \qquad \gamma_d(\mathcal E_d)\leq e^{-4d^2}.
\]
\end{lemma}

\begin{proof}
For the first bound use the witness in \eqref{eq:matrixblocks}, before
adding the noise. Conditionally on $A^{(d)}$, its projections are fixed
and independent of $G^{(d)}$. Put $r_l=\rank(P_l)$. The block-diagonal
Hermitian subspace has real dimension $\sum_l r_l^2$, so the GUE
covariance $1/d$ and the normalized trace give
\[
 \E\big[\|\sqrt t\D_P G^{(d)}_{[p]}\|_{2,d}^2\mid A^{(d)}\big]
 =tp\sum_l(\rank(P_l)/d)^2\leq2t.
\]
Markov's inequality at squared threshold $p/16$ bounds failure by
$32t/p<1/8$. On success the projected total norm exceeds
$\sqrt{p/2}-\sqrt p/4>\sqrt p/4$.
Independence and centering of the noise also give
\[
 \E\|A^{(d)}_{[p]}+\sqrt t G^{(d)}_{[p]}\|_{2,d}^2
 =\E\|A^{(d)}_{[p]}\|_{2,d}^2+tp\longrightarrow vp.
\]
Markov's inequality at the squared norm $16vp$ bounds the failure of
the norm condition in \eqref{eq:event} by $1/16+o(1)<1/8$.
The union bound shows that both conditions hold with probability at
least $3/4$, proving the weaker claimed bound.

For the Gaussian bound enumerate at most $(d+1)^p$ rank lists and choose
an operator-norm $\delta$-net of the unitary group, with
$\delta=1/(64\sqrt v)$. A maximal $\delta$-separated subset of $U(d)$
has unitary centers and is a $\delta$-net. Its ambient operator-norm
balls of radius $\delta/2$ are disjoint and lie in the ball of radius
$1+\delta/2$ in the real $2d^2$-dimensional space $M_d(\C)$. Volume
comparison gives cardinality at most $(1+2/\delta)^{2d^2}$.
Every decomposition of a fixed rank list is a unitary conjugate of
its standard coordinate blocks $P^0$, including lists with zero ranks.
For unitaries $U,V$ with $\|U-V\|\leq\delta$,
\[
 \|U^*D_{[p]}U-V^*D_{[p]}V\|_{2,d}
 \leq2\delta\|D_{[p]}\|_{2,d}.
\]
Fixed-block compression is contractive. Thus replacing a witnessing
unitary by a net point loses at most
$2\delta\cdot4\sqrt{vp}=\sqrt p/8$ on \eqref{eq:event}.
Hence some rank list satisfying the bounds and some net unitary $U$
must satisfy
\[
 \|\D_{P^0}(U^*D_{[p]}U)\|_{2,d}\geq\sqrt p/8.
\]
For a fixed list and unitary, unitary invariance preserves the independent
variance-$v$ GUE law of the first $p$ coordinates. Their projected
Hilbert--Schmidt vector has
$q=p\sum_l\rank(P_l^0)^2\leq2d^2$ real coordinates of variance $v/d$.
The further division by $d$ in the normalized squared norm gives the
law $(v/d^2)Q$, where $Q$ is chi-squared of degree $q$.
Since $\E e^{Q/4}=2^{q/2}$, its probability of reaching the squared-norm
threshold $p/64$, equivalently $Q\geq pd^2/(64v)$, is at most
\[
 \exp\left[d^2\left(\log2-\frac{p}{256v}\right)\right].
\]
The full union bound consequently has logarithm at most
\[
 p\log(d+1)+d^2\left[2\log(1+128\sqrt v)+\log2-\frac{p}{256v}\right].
\]
Since $v<2+t$, Equation \eqref{eq:parameters} makes the bracket less
than $-8$. The first term is $o(d^2)$ because $p$ is fixed. This proves
the bound $e^{-4d^2}$ for large $d$.
\end{proof}

\begin{proof}[Proof of Theorem~\ref{thm:main}]
Retain the same cutoff, dimension subsequence, and ensembles throughout.
Put $\alpha=\mu_d(\mathcal E_d)$ and
$\beta=\gamma_d(\mathcal E_d)$. Binary data processing for relative
entropy and the bound on binary entropy give
\[
 \KL(\mu_d\Vert\gamma_d)
 \geq\alpha\log\frac\alpha\beta
 +(1-\alpha)\log\frac{1-\alpha}{1-\beta}
 \geq\alpha\log\frac1\beta-\log2.
\]
Thus Lemma~\ref{lem:eventprob} implies
$\liminf_d d^{-2}\KL(\mu_d\Vert\gamma_d)\geq2$.
We retain only one unit of this normalized deficit in the comparison below.
At positive time the bounded-input Gaussian convolution has finite
entropy and finite second moments, and its relative entropy with respect
to $\gamma_d$ is finite. The Gaussian reference has density
\[
 g_d(D)=\prod_{i=1}^n
 \left[
 \left(\frac{d}{2\pi v_i(h+t)}\right)^{d^2/2}
 \exp\left(-\frac{d^2\tr_d(D_i^2)}{2v_i(h+t)}\right)
 \right].
\]
Using $-\int f\log f=-\int f\log g_d-\KL(f\Vert g_d)$ and the
normalization of $H_d$ gives the exact identity
\begin{align*}
 H_d(Z^{(d)}(t))
 &=\frac12\sum_i\left[
 \log(2\pi v_i(h+t))+
 \frac{\E\tr_d(Z_i^{(d)}(t)^2)}{v_i(h+t)}\right]
 -d^{-2}\KL(\mu_d\Vert\gamma_d).
\end{align*}
The expected second moments equal those of $A_i^{(d)}$ plus $t$, and
converge to $v_i(h+t)$. In particular,
\[
 \limsup_d H_d(Z^{(d)}(t))\leq L(h+t)-1.
\]
This upper bound concerns exactly the terminal law
$Z^{(d)}(t)=A^{(d)}+\sqrt t G^{(d)}$ in Lemma~\ref{lem:flow}.
Combining the two bounds on its entropy and then
\eqref{eq:largeentropy}--\eqref{eq:freeflow} gives
\begin{align*}
 \chi_R(X)
 &\leq L(h+t)-1-\frac12\int_0^t\Phi^*(V_{h+s})\,ds\\
 &\leq\chi^*(V_{h+t})-\tfrac12
       -\frac12\int_0^t\Phi^*(V_{h+s})\,ds\\
 &=\chi^*(V_h)-\tfrac12.
\end{align*}
This holds for every fixed cutoff with $\chi_R(X)>-\infty$; for the
remaining cutoffs it is automatic. Taking the supremum over cutoffs
preserves the bound. Lemma~\ref{lem:finitechi} gives $\chi(X)>-\infty$,
and \eqref{eq:entropydeficit} shows that $\chi^*(X)$ is a finite real number.
The bounded matrix realization in Section~\ref{sec:signal} supplies
the asserted faithful normal tracial state and bounded self-adjoint tuple.
\end{proof}
